% Supplementary C — Conference paper §III (Experimental setup)
% Full dataset descriptors, tokenisation choices, and per-scene
% sampling traceability.

\documentclass[11pt,a4paper]{article}
% Shared preamble for all supplementary chapter documents.
% Used by each supplementary/conference/suppl_*.tex and
% supplementary/journal/suppl_*.tex via % Shared preamble for all supplementary chapter documents.
% Used by each supplementary/conference/suppl_*.tex and
% supplementary/journal/suppl_*.tex via % Shared preamble for all supplementary chapter documents.
% Used by each supplementary/conference/suppl_*.tex and
% supplementary/journal/suppl_*.tex via \input{../preamble.tex}.

\usepackage[utf8]{inputenc}
\usepackage[T1]{fontenc}
\usepackage[a4paper,margin=2.2cm]{geometry}
\usepackage{amsmath,amssymb,amsfonts,amsthm}
\usepackage{mathtools}
\usepackage{graphicx}
\usepackage{booktabs}
\usepackage{array}
\usepackage{multirow}
\usepackage{longtable}
\usepackage{algorithm}
\usepackage{algpseudocode}
\usepackage{xcolor}
\usepackage{url}
\usepackage{cite}
\usepackage[hidelinks]{hyperref}

\graphicspath{{../../figures/}}

\renewcommand{\thesection}{\Alph{section}.\arabic{section}}
\renewcommand{\thesubsection}{\thesection.\arabic{subsection}}

\theoremstyle{plain}
\newtheorem{theorem}{Theorem}[section]
\newtheorem{lemma}[theorem]{Lemma}
\newtheorem{proposition}[theorem]{Proposition}
\theoremstyle{definition}
\newtheorem{definition}[theorem]{Definition}
\theoremstyle{remark}
\newtheorem{remark}[theorem]{Remark}

\setcounter{secnumdepth}{3}
\setcounter{tocdepth}{2}

\newcommand{\R}{\mathbb{R}}
\newcommand{\E}{\mathbb{E}}
\newcommand{\KL}[2]{\mathrm{KL}\!\left(#1 \,\middle\|\, #2\right)}
\newcommand{\ari}{\mathrm{ARI}}
\newcommand{\nmi}{\mathrm{NMI}}
\newcommand{\softmax}{\mathrm{softmax}}
.

\usepackage[utf8]{inputenc}
\usepackage[T1]{fontenc}
\usepackage[a4paper,margin=2.2cm]{geometry}
\usepackage{amsmath,amssymb,amsfonts,amsthm}
\usepackage{mathtools}
\usepackage{graphicx}
\usepackage{booktabs}
\usepackage{array}
\usepackage{multirow}
\usepackage{longtable}
\usepackage{algorithm}
\usepackage{algpseudocode}
\usepackage{xcolor}
\usepackage{url}
\usepackage{cite}
\usepackage[hidelinks]{hyperref}

\graphicspath{{../../figures/}}

\renewcommand{\thesection}{\Alph{section}.\arabic{section}}
\renewcommand{\thesubsection}{\thesection.\arabic{subsection}}

\theoremstyle{plain}
\newtheorem{theorem}{Theorem}[section]
\newtheorem{lemma}[theorem]{Lemma}
\newtheorem{proposition}[theorem]{Proposition}
\theoremstyle{definition}
\newtheorem{definition}[theorem]{Definition}
\theoremstyle{remark}
\newtheorem{remark}[theorem]{Remark}

\setcounter{secnumdepth}{3}
\setcounter{tocdepth}{2}

\newcommand{\R}{\mathbb{R}}
\newcommand{\E}{\mathbb{E}}
\newcommand{\KL}[2]{\mathrm{KL}\!\left(#1 \,\middle\|\, #2\right)}
\newcommand{\ari}{\mathrm{ARI}}
\newcommand{\nmi}{\mathrm{NMI}}
\newcommand{\softmax}{\mathrm{softmax}}
.

\usepackage[utf8]{inputenc}
\usepackage[T1]{fontenc}
\usepackage[a4paper,margin=2.2cm]{geometry}
\usepackage{amsmath,amssymb,amsfonts,amsthm}
\usepackage{mathtools}
\usepackage{graphicx}
\usepackage{booktabs}
\usepackage{array}
\usepackage{multirow}
\usepackage{longtable}
\usepackage{algorithm}
\usepackage{algpseudocode}
\usepackage{xcolor}
\usepackage{url}
\usepackage{cite}
\usepackage[hidelinks]{hyperref}

\graphicspath{{../../figures/}}

\renewcommand{\thesection}{\Alph{section}.\arabic{section}}
\renewcommand{\thesubsection}{\thesection.\arabic{subsection}}

\theoremstyle{plain}
\newtheorem{theorem}{Theorem}[section]
\newtheorem{lemma}[theorem]{Lemma}
\newtheorem{proposition}[theorem]{Proposition}
\theoremstyle{definition}
\newtheorem{definition}[theorem]{Definition}
\theoremstyle{remark}
\newtheorem{remark}[theorem]{Remark}

\setcounter{secnumdepth}{3}
\setcounter{tocdepth}{2}

\newcommand{\R}{\mathbb{R}}
\newcommand{\E}{\mathbb{E}}
\newcommand{\KL}[2]{\mathrm{KL}\!\left(#1 \,\middle\|\, #2\right)}
\newcommand{\ari}{\mathrm{ARI}}
\newcommand{\nmi}{\mathrm{NMI}}
\newcommand{\softmax}{\mathrm{softmax}}


\title{Supplementary C --- Dataset descriptors and tokenisation\\
{\large Companion to: \emph{A Band-Mask Robustness Diagnostic for LDA on HSI}}}
\author{Felipe Santib\'a\~nez-Leal}
\date{2026}

\begin{document}
\maketitle

\section{Scope}
This supplement expands Section~III of the conference paper. We
document per-scene the sensor history, band selection, ground-truth
class catalogue, and the sampling decisions that produced the
canonical 220-per-class corpora used in the band-mask diagnostic.

\section{Per-scene sensor and corpus traceability}

\begin{longtable}{l p{2.6cm} c c c p{4.0cm}}
\caption{Labelled HSI scenes used in the diagnostic. $B$ is the band
count after the canonical correction (uncorrected bands removed per
the UPV/EHU collection's official policy). $D$ is the document count
that enters the canonical fit (= samples per class × number of
admissible classes, with classes having fewer than $100$ available
pixels rejected). $K$ is the canonical topic count.}
\label{tab:scenes-full}\\
\toprule
Scene & Sensor / platform & $B$ & $D$ & $K$ & Source URL  \\
\midrule
Indian Pines       & AVIRIS / NASA      & 200 & 2812 & 12 & ehu.eus/ccwintco/index.php?title=Hyperspectral\_Remote\_Sensing\_Scenes \\
Salinas            & AVIRIS / NASA      & 204 & 3520 & 12 & ehu.eus / Salinas \\
Salinas-A          & AVIRIS / NASA      & 204 & 1320 & 6  & ehu.eus / Salinas-A (subset)  \\
Pavia U            & ROSIS-03 / DLR     & 103 & 1980 & 9  & ehu.eus / Pavia University \\
KSC                & AVIRIS / NASA      & 176 & 2686 & 12 & ehu.eus / Kennedy Space Center \\
Botswana           & EO-1 Hyperion / NASA & 145 & 2775 & 12 & ehu.eus / Botswana \\
\bottomrule
\end{longtable}

\subsection{Indian Pines}
\textbf{Sensor.} Airborne Visible/Infrared Imaging Spectrometer
(AVIRIS), 224 raw bands at $\approx 10$~nm resolution over
400--2500~nm. The canonical correction removes 24 noisy and
atmospheric-water bands, leaving $B = 200$ usable bands.
\textbf{Scene.} $145 \times 145$ pixels acquired over the
agricultural test site at Purdue University, Indiana, June 1992
\cite{baumgardner1992indianpines}. \textbf{Classes.} 16 ground-truth
classes covering corn-notill, corn-mintill, corn, grass-pasture,
grass-trees, hay-windrowed, oats, soybean-notill, soybean-mintill,
soybean-clean, wheat, woods, buildings-grass-trees-drives,
stone-steel-towers, alfalfa, grass-pasture-mowed. Some classes have
fewer than $100$ labelled pixels and are excluded from the canonical
sample.
\textbf{Sampling.} For each admissible class, the canonical fit draws
$220$ pixels without replacement using NumPy's
\texttt{default\_rng(42)}.

\subsection{Salinas}
\textbf{Sensor.} AVIRIS, same instrument family as Indian Pines.
\textbf{Scene.} $512 \times 217$ pixels acquired over Salinas Valley,
California~\cite{aviris_salinas}, $B = 204$ usable bands after the
canonical correction.
\textbf{Classes.} 16 classes: brocoli-green-weeds-1/2, fallow,
fallow-rough-plow, fallow-smooth, stubble, celery, grapes-untrained,
soil-vinyard-develop, corn-senesced-green-weeds,
lettuce-romaine-{4,5,6,7}wk, vinyard-untrained, vinyard-vertical-trellis.

\subsection{Salinas-A}
\textbf{Sensor + classes.} A subset of Salinas comprising 6 of the
classes on a $83 \times 86$ window. Sensor history identical to
Salinas. We treat it as an independent benchmark because its smaller
spatial footprint and fewer classes make the canonical fit
substantially easier than the parent scene.

\subsection{Pavia University}
\textbf{Sensor.} Reflective Optics System Imaging Spectrometer
(ROSIS-03), DLR. 115 raw bands over 430--860~nm; the canonical
correction removes 12 noisy bands leaving $B = 103$, all in the VNIR
range. The scene is consequently \emph{auto-skipped} for the SWIR
mask in the diagnostic ($V^{(m)} = 0$ bands kept).
\textbf{Scene.} $610 \times 340$ pixels over the University of Pavia,
Italy~\cite{rosis_pavia}.
\textbf{Classes.} 9 classes: asphalt, meadows, gravel, trees,
painted-metal-sheets, bare-soil, bitumen, self-blocking-bricks,
shadows.

\subsection{Kennedy Space Center (KSC)}
\textbf{Sensor.} AVIRIS-Old (AVIRIS recording before 1996).
\textbf{Scene.} $512 \times 614$ pixels acquired over the Kennedy
Space Center, Florida, March 1996. $B = 176$ usable bands after the
canonical correction~\cite{aviris_ksc}.
\textbf{Classes.} 13 classes covering wetland vegetation
(salt-marsh, mud-flat, water, swamp, etc.).
\textbf{Note.} KSC is the labelled scene most frequently observed to
yield a band-fragile canonical basis, an observation surfaced by the
diagnostic; the supplementary E document discusses possible
mechanisms in more depth.

\subsection{Botswana}
\textbf{Sensor.} Earth Observing 1 (EO-1) Hyperion.
\textbf{Scene.} $1476 \times 256$ pixels acquired over the Okavango
Delta, Botswana, May 2001~\cite{eo1_botswana}. $B = 145$ usable bands
after the canonical correction.
\textbf{Classes.} 14 classes including water, hippo-grass,
floodplain-grasses-1/2, reeds, riparian, firescar-2, island-interior,
acacia-woodlands, acacia-shrublands, acacia-grasslands,
short-mopane, mixed-mopane, exposed-soils.

\section{Tokenisation reproducibility note}

The min-max normalisation of \S2.1 of supplementary B is computed
\emph{within each pixel} and not across the corpus, so the
tokenisation of pixel $d$ depends only on the spectrum of pixel $d$.
This is by design: across-corpus normalisation would couple the
tokenisation to the sample selection, and the band-mask diagnostic
would then conflate spectral-window restriction with corpus-extent
sensitivity. The trade-off is that pixels with very flat spectra
(e.g.\ heavy shadow) tokenise to a uniform count vector, which the
LDA fit downweights via the small $\alpha$ prior.

\section{HIDSAG mineral subsets}

The five HIDSAG subsets evaluated alongside the labelled scenes are
documented in detail in the journal supplement E. Briefly: HIDSAG is
a curated collection of point-spectrum measurements over
mineralogical samples acquired with a SWIR-only sensor; the five
subsets (GEOMET, MINERAL1, MINERAL2, GEOCHEM, PORPHYRY) collectively
cover 19 target labels distributed across multiple samples. The
band-mask diagnostic on HIDSAG uses \emph{between-covariate-variance}
band ranking in place of the labelled-scene Fisher discriminant, as
HIDSAG has no per-pixel ground truth.

\bibliographystyle{IEEEtran}
\bibliography{../../bibliography/refs}

\end{document}
